\subsection{Counting maps with labelled cells}\label{sec:pinned-count}
An action on cells has both an image and a kernel, and the two play
different roles in a counting problem. Let an oriented map have finite
dart set $E$, vertex rotation $\rho$, and edge involution $\alpha$;
we use $\phi=\rho^{-1}\alpha$ for its face permutation. Fixed points
of $\alpha$ are allowed and describe half-edges. Connectedness means
transitivity of $\langle\rho,\alpha\rangle$.

The classical covering correspondence identifies transitive permutation
data with Belyi pairs and cycle lengths with ramification indices
\cite[Lemma~I.1, p.~2, and Theorem~I.5, p.~8]{Schneps}. A permutation
commuting with both generators gives a unique covering transformation
by path lifting. The local power models extend it over the branch
points. Thus the Belyi automorphism group is their simultaneous
centralizer. For degree $d$ the genus is determined by
\[
 \#\operatorname{cyc}(\rho)+\#\operatorname{cyc}(\alpha)
 +\#\operatorname{cyc}(\phi)=d+2-2g.
\]
In the clean case this becomes the usual vertices-minus-edges-plus-faces
formula, since $d$ is twice the number of ordinary edges.

Fix a type $T$ consisting of graded sets of vertices, edges and faces,
where the grade is the degree. Write
$\gamma_T=|\operatorname{Aut}(T)|$; if the counts in degree $j$ are
$a_j,b_j,c_j$, then $\gamma_T=\prod_j a_j!b_j!c_j!$.
A pinning labels the graded cells of a map by $T$. For an unlabelled
map $C$, let $A_C$ be its full combinatorial group, and let $K_C,I_C$
be the kernel and image of its action on those graded cells.

\begin{proposition}\label{prop:pinning-count}
If $\mathcal M(T)$ is the finite set of unpinned isomorphism classes
of type $T$, with any specified connectedness or genus restrictions,
then the number of pinned isomorphism classes is
\[
 N_T=\sum_{[C]\in\mathcal M(T)}\frac{\gamma_T}{|I_C|}
 =\gamma_T\sum_{[C]\in\mathcal M(T)}\frac{|K_C|}{|A_C|}.
\]
Each pinning of $C$ has automorphism group $K_C$. Equality with
$|\mathcal M(T)|\gamma_T$ holds exactly when every $I_C$ is trivial;
pinned rigidity instead means every $K_C$ is trivial.
\end{proposition}
\begin{proof}
For fixed $C$, its pinning set is a torsor under $\operatorname{Aut}(T)$.
The action of $A_C$ factors through $I_C$, whose action on this torsor
is free. Hence every orbit has size $|I_C|$, while every stabilizer
in $A_C$ is $K_C$. Sum the orbit counts over distinct unpinned maps.
Each summand is at most $\gamma_T$, with equality precisely when
$|I_C|=1$, proving also the asserted criterion. The argument includes
disconnected maps and the empty class set. If the empty map is allowed,
it has one pinning of empty type; a nonempty clean map has even dart
number.
\end{proof}

For the smallest counterexample to identifying these conditions, take
$E=\Z/4$, $\rho(x)=x+1$ and $\alpha(x)=x+2$.
Here $\phi=\rho$, so there is one vertex, two edges and one face,
on a torus. The centralizer is the translation group $C_4$; its even
translations fix the two edges and its odd translations interchange
them. Thus $K_C=C_2$, $I_C=C_2$ and $\gamma_T=2$.
There is only one unpinned map of this type: of the three pairings
of four cyclically ordered darts, only the opposite pairing gives
one face. The formula therefore gives one pinned class. The kernel
classification in the proof of Theorem~\ref{thm:dessins} identifies
this as the genus-one member of the complete one-vertex one-face
family, and shows that any half-edge forces $K_C=1$.

\subsection{Functions on fixed embedded cells and noncompact maps}
The group in Theorem~\ref{thm:dessins} preserves the reduced map,
whose white points have been suppressed. It differs from the group
commuting with its Belyi function. This distinction also matters for
noncompact dessins: Cañas, Hidalgo, Turiel and Viruel realize every
countable group as a dessin automorphism group
\cite[Theorems~1--2, p.~3; Section~2.3, p.~4]{Canas}.
Their automorphisms may permute cells. The group considered here
fixes every cell separately.

\begin{proposition}\label{prop:fixed-cell-torsor}
Fix a connected compact Riemann surface $X$ with a clean dessin's embedded
reduced cells, and let $H$ be the subgroup of its holomorphic
cellwise group fixing every dart. The clean Belyi functions having
exactly these cells form a free transitive $H$-set under precomposition.
Thus nonuniqueness occurs exactly for the one-loop sphere and the
oranges of every order. Marking every face centre removes these
continuous groups. Marking every white point removes them for
oranges but leaves positive dilations for the one-loop sphere.
\end{proposition}
\begin{proof}
Two such functions give identical oriented dart data. The covering
isomorphism specified by the identity of these data extends
holomorphically across each branch point, by the local power-map
argument above. It fixes each dart and intertwines the functions,
so belongs to $H$. Conversely precomposition by $H$ preserves the
embedded cells. A stabilizer would be a Belyi automorphism fixing a
dart, and connectedness then makes it the identity. This proves the
torsor assertion. The compact classification identifies $H$ with
the positive affine group or positive dilations in the two spherical
families, and with the trivial group in positive genus. The finite
positive-genus involution already preserves the Belyi function.

For $1/(1+z^2)$ the white point is $0$ and the two face centres are
$i,-i$. An affine map $az+b$, $a>0,b\in\R$, fixes $0$ precisely when
$b=0$, and fixes $i$ only when also $a=1$. In an orange every white
point and face centre is finite and nonzero, so fixing any one of them
forces a positive dilation to have scale one.
\end{proof}
The torus involution reverses the two loops of its one-vertex
one-face structure and acts as $-1$ on $H_1(-;\Z)$; it is therefore
not isotopic to the identity. Conversely a M\"obius transformation is
isotopic to the identity on the sphere but can move its dessin.
Cellwise preservation and isotopy to the identity are different conditions.

\begin{proposition}\label{prop:noncompact-cell-rigidity}
Let a connected noncompact Riemann surface carry a locally finite
clean holomorphic cellular map, with compact edge closures and
ordinary disk faces attached along their finite edge cycles.
Every cellwise holomorphic automorphism is the identity.
\end{proposition}
\begin{proof}
The transitive dart action has a centralizer acting freely on darts.
If a nonidentity permutation in it fixes each edge, it must reverse
every edge, hence equal $\alpha$. Fixing every vertex then forces
one vertex rotation cycle. Local finiteness makes that cycle finite,
so there are finitely many darts, edges and face cycles. Their
ordinary closed-cell attachments make the surface compact, a
contradiction. Thus any cellwise automorphism fixes all darts.

Its derivative at a black vertex is positive real because it
preserves each outgoing analytic ray. Uniformization
\cite[Theorem~4.2.1 and the quotient classification, pp.~152--153]{Wendl}
reduces to a disk or plane universal cover. In the disk case lift
the map to fix a chosen lift of the vertex. After moving it to zero,
Schwarz's lemma makes the lift a rotation; its positive derivative
forces the identity. For a plane cover with a nontrivial deck group,
noncompactness leaves a rank-one translation lattice. The lift fixing
a chosen vertex is affine, and its positive multiplier preserves
this lattice, so is one. Finally, on $\C$ a nonidentity positive
dilation about the fixed vertex has a two-sided orbit unbounded in
one direction. An interior point of an invariant compact edge
closure cannot have such an orbit. The map is therefore the identity
in every case.
\end{proof}
The ordinary-cell hypothesis excludes asymptotic tracts; an arbitrary
holomorphic map to the sphere need not have the required cellular structure.

\section{Schneps's sextic example and its normalization}\label{sec:sextic}
This section makes the known three-member sextic orbit explicit in the
normalization suggested by the manuscript. Schneps's cubic in
\cite[Section~IV, Example~3, pp.~27--28 of the author version]{Schneps}
is $25t^3-6t^2-6t-2$; setting $c=2t$ gives the cubic below.
The cubic and its three-member orbit are credited to that example.
The factorization here also separates the collision parameters from
the required passport.

Set
\begin{align*}
 h_c(z)&=z^3(z-1)^2(z-c),&
 G_c(z)&=z^2-\frac{5c+4}{6}z+\frac c2,\\
 \delta(c)&=25c^2-32c+16,&p(c)&=25c^3-12c^2-24c-16,\\
 P(c)&=-\frac{\delta(c)p(c)}{7776},&
 Q(c)&=\frac{c(5c-8)(25c^3-6c^2+8)}{2592}.
\end{align*}
Differentiation gives $h_c'=6z^2(z-1)G_c$. Division by $G_c$ gives
remainder $P(c)z+Q(c)$. To verify it directly, put
$\alpha=(5c+4)/6$, $\beta=-c/2$, so $z^2=\alpha z+\beta$
in the quotient. Reduce the six powers in $h_c$ using this relation;
the resulting coefficients are
\begin{align*}
 7776P&=-625c^5+1100c^4-184c^3-176c^2-128c+256,\\
 2592Q&=125c^5-230c^4+48c^3+40c^2-64c.
\end{align*}
Factoring them gives the displayed formulas. If $G_c$ has distinct
roots $u,v$, their critical values agree precisely when
$P(c)(u-v)=0$, hence, after excluding $\delta=0$, when $p(c)=0$.
The raw remainder coefficient has degree five; the effective
noncollision equation has degree three.

\begin{proposition}\label{prop:sextic-complete}
The three roots of $p$ give exactly the Belyi pairs with ordered
passport $[3,2,1;2,2,1,1;6]$ through
\[
 \beta_c(z)=\frac{1500z^3(z-1)^2(z-c)}{32-52c-c^2}.
\]
They have trivial Belyi automorphism groups. Writing
$r_j=\sqrt[3]{2}\,\zeta_3^j$ with positive real $\sqrt[3]{2}$,
the parameters are
$c_j=(4+18r_j+6r_j^2)/25$. Their fields of moduli and definition are
$\Q(c_j)=\Q(r_j)$; their splitting field is
$\Q(\sqrt[3]{2},\zeta_3)$, with Galois group $S_3$ and quadratic
subfield $\Q(\sqrt{-3})$.
\end{proposition}
\begin{proof}
The values of $p$ modulo $7$ at $0,\ldots,6$ are
$5,1,4,3,1,1,6$. Thus $p$ is irreducible over $\Q$.
Its discriminant is $-2^8 3^9$, so it has one real and two nonreal
roots, all distinct. Modulo $p$,
\[
 Q=\frac{32-52c-c^2}{1500},\qquad
 h_c-Q=G_c^2H_c,
 \quad H_c=z^2+\frac{2c-2}{3}z+\frac{5c^2-4c-4}{12}.
\]
For an explicit identity before reduction, the difference
$h_c-Q-G_c^2H_c$ is
\[
 \frac{p(c)}{2592}
 (-5c^2-4c-30cz^2+36cz+24z^3-24z^2).
\]
None of $c,1-c,Q,\delta$ vanishes at a root of the irreducible cubic.
The discriminant of $H_c$ is $(16+4c-11c^2)/9$, also nonzero.
The resultant of $G_c,H_c$, reduced modulo $p$, is
$(c-1)(5c^2+2c+2)/6$, so these factors have no common root.
Their zeros give exactly two double and two simple white vertices.
The derivative has no other finite critical points, and infinity is
a pole of order six. This proves the passport.

Conversely that passport has total ramification $3+2+5=10$;
Riemann--Hurwitz makes its source a sphere. Put its unique pole at
infinity, its triple black point at zero and its double black point
at one. The normalized function is uniquely $\lambda h_c$.
Its residual critical points are distinct and have value one, forcing
$p(c)=0$ and $\lambda=Q(c)^{-1}$. Every pair is therefore on the list.
An isomorphism fixes the three distinguished points and hence is the
identity; it also forces equality of the parameters.

Substitute $c=(4+18r+6r^2)/25$ and reduce by $r^3-2$ to obtain
$p(c)=0$. Both $\Q(c)$ and $\Q(r)$ have degree three and the former
is contained in the latter, so they coincide. The three substitutions
give all conjugates. The irreducible cubic's nonsquare discriminant
makes the Galois group $S_3$; its square root generates
$\Q(\sqrt{-3})$. Finally, Galois conjugation sends $\beta_c$ to
$\beta_{\sigma(c)}$. The unique normalization shows that these pairs
are isomorphic exactly when $\sigma(c)=c$. This identifies the field
of moduli directly, and the formula gives a model over it.
\end{proof}

The corresponding monodromy representatives, composing right to left,
are
\[
\begin{array}{c|c|c}
 &\alpha&\phi=\rho^{-1}\alpha\\\hline
 T_1&(14)(36)&(154362)\\
 T_0&(14)(56)&(156432)\\
 T_2&(14)(26)&(154326)
\end{array}
\qquad\rho=(123)(45)(6).
\]
Each face permutation is a six-cycle, so each map is connected and
has genus zero. The centralizer of $\rho$ is $C_3\times C_2$.
In $T_0$ the singleton dart $6$ is paired with a dart of the two-cycle
of $\rho$, whereas in $T_1,T_2$ it is paired with a dart of the
three-cycle. The first distinction is invariant. To carry $T_1$ to
$T_2$, a centralizing permutation would have to carry $3$ to $2$;
its unique rotation on $(123)$ then carries the other paired dart
$1$ to $3$, not to $1$, so it cannot carry $(14)$ to $(14)$.
Thus the three classes are distinct, and the geometric completeness
proof shows they exhaust the passport. Conjugation by $(23)$ reverses
$\rho$ and exchanges $T_1,T_2$, while fixing $T_0$.
The real parameter therefore gives $T_0$, and the nonreal parameters
give the mirror pair. A separate exact enumeration finds all eighteen
one-face involutions in three centralizer orbits of six.

For completeness, the roots of $\delta$ are
$16/25\pm12i/25$. The repeated residual critical point
$u=(5c+4)/12$ is distinct from $0,1,c$, and
$Q\equiv-2(139c-82)/9375\pmod\delta$ is nonzero.
The normalized polynomial then has white passport $[3,1,1,1]$,
because its derivative has order two at $u$ and no other residual
critical point. These parameters must be excluded from the preceding
proposition. The ancillary coefficient identities are
\[
 (\alpha-2)(\alpha-c)+\beta+1
 =-\frac{5c^2-10c-4}{36},\qquad
 \alpha(\alpha^2+2\beta)
 =\frac{(5c+4)(25c^2+4c+16)}{216}.
\]
The two double white points are the roots of $G_c$; their multiplicity-two factor in $f-1$ is $G_c^2$.
More generally, logarithmic differentiation of coprime factorizations
gives
\[
 \frac{f'}f=\sum_i\frac{p_i}{z-a_i}-\sum_j\frac{q_j}{z-b_j}
 \quad\text{for }f=\lambda
 \frac{\prod_i(z-a_i)^{p_i}}{\prod_j(z-b_j)^{q_j}}.
\]
